\uuid{5hei}
\titre{Dérivée en un point}
\auteur{bodin ; hanani ; mzari}
\organisation{exo7}
\licence{CC BY-NC-SA 4.0}
\datecreate{2020-01-20}
\niveau{L1}
\module{Analyse}
\chapitre{Dérivabilité des fonctions réelles}
\sousChapitre{Autre}
\difficulte{1}
\theme{Dérivées, dérivée en un point}
\exo7qcmid{lille-1.derivees.003}

\contenu{
\begin{question}
\qclassification{124.00}
\qtags{motcle=dérivée en un point}
Soit $\displaystyle f(x)=|x-1|$. On note $f'_d(a)$ (resp. $f'_g(a)$) pour désigner la dérivée à droite (resp. à gauche) en $a$. Quelles sont les bonnes réponses ?
\begin{answers}
    \bad{$f'_d(1)=1$ et $f'_g(1)=1$}
    \bad{$f$ est dérivable en $1$ et $f'(1)=1$.}
    \good{$f$ est dérivable en $0$ et $f'(0)=-1$.}
    \good{$f$ n'est pas dérivable en $1$ car $f'_d(1)=1$ et $f'_g(1)=-1$.}
\end{answers}
\begin{explanations}
Par définition, on a :
$$f(x)=\left\{ \begin{array}{ll}x-1&\mbox{si }x\geq 1\\ 1-x&\mbox{si }x\leq 1.
\end{array}\right.$$
Donc, $f$ est dérivable sur $\Rr\setminus\{1\}$, et
$$f'(x)=\left\{ \begin{array}{ll}1&\mbox{si }x> 1\\ -1&\mbox{si }x< 1.
\end{array}\right.$$
En particulier, $f'(0)=-1$. Par contre $f$ n'est pas dérivable en $1$ car
$$\lim _{x\to 1^+}\frac{f(x)-f(1)}{x-1}=1 \mbox{ et }\lim _{x\to 1^-}\frac{f(x)-f(1)}{x-1}=-1.$$
\end{explanations}
\end{question}
}
